\documentclass[10pt]{article}
\usepackage[T1]{fontenc}
\usepackage[utf8]{inputenc}
\usepackage[french]{babel}
\usepackage{amsmath,amssymb,mathrsfs,mathpazo}
\usepackage{geometry}
\geometry{a4paper,margin=2cm}
\usepackage{graphicx,array,tabularx,tabularray,colortbl,multirow,multicol,booktabs,enumitem,pifont,ulem,fancybox,diagbox,cancel}
\usepackage{pstricks,pstricks-add,pst-plot,pst-tree,pst-node,pst-text,pst-func,pst-math,pst-bspline,pst-3dplot,pst-eucl,pst-solides3d}
\let\pscalculate\undefined
\usepackage{pst-bezier}
\usepackage{tikz,pgfplots,tkz-tab}
\usetikzlibrary{arrows,arrows.meta,patterns,calc,positioning,trees,decorations.markings,shapes.geometric,shapes.misc,shapes.arrows,intersections,decorations.pathmorphing,decorations.pathreplacing,backgrounds,shadows,babel}
\pgfplotsset{compat=1.18}
\usepgfplotslibrary{fillbetween}
\usepackage[np]{numprint}
\usepackage{listings,makecell,pas-tableur,twoopt}
\let\barre\undefined
\usepackage[pstricks]{tablvar}
\definecolor{darkgreen}{RGB}{0,100,0}
\providecommand{\red}{\color{red}}
\providecommand{\blue}{\color{blue}}
\providecommand{\green}{\color{green}}
\providecommand{\black}{\color{black}}
\providecommand{\white}{\color{white}}
\providecommand{\gray}{\color{gray}}
\providecommand{\cyan}{\color{cyan}}
\providecommand{\magenta}{\color{magenta}}
\providecommand{\yellow}{\color{yellow}}
\providecommand{\violet}{\color{violet}}
\providecommand{\orange}{\color{orange}}
\providecommand{\euro}{\texteuro}
\providecommand{\EUR}{\texteuro}
\providecommand{\ieme}{\textsuperscript{e}}
\providecommand{\degres}{\ensuremath{^\circ}}
\providecommand{\degree}{\ensuremath{^\circ}}
\providecommand{\vv}[1]{\overrightarrow{#1}}
\providecommand{\Coord}[1]{#1}
\providecommand{\decosix}{}
\providecommand{\starredbullet}{\ensuremath{\bullet}}
\pagestyle{empty}
\setlength{\parindent}{0pt}
\usepackage{adjustbox}\DeclareUnicodeCharacter{2236}{:}\DeclareUnicodeCharacter{2190}{\ensuremath{\leftarrow}}
\providecommand{\R}{\mathbb{R}}
\providecommand{\N}{\mathbb{N}}
\providecommand{\Z}{\mathbb{Z}}
\providecommand{\Q}{\mathbb{Q}}
\providecommand{\D}{\mathbb{D}}
\providecommand{\C}{\mathbb{C}}
\providecommand{\e}{{\mathrm{e}}}
\providecommand{\vect}[1]{\overrightarrow{#1}}
\providecommand{\vectt}[1]{\overrightarrow{\mathrm{#1}}}
\providecommand{\vvt}[1]{\overrightarrow{\mathrm{#1}}}
\providecommand{\np}[2][]{{#2~#1}}
\providecommand{\vv}[1]{{\overrightarrow{#1}}}
\renewcommand{\d}{{\mathrm{d}}}
\providecommand{\degres}{{^\circ}}
\providecommand{\degre}{{^\circ}}
\providecommand{\blue}{\color{blue}}
\providecommand{\red}{\color{red}}
\providecommand{\enskip}{\quad}
\providecommand{\llbracket}{\lbrack\!\lbrack}
\providecommand{\rrbracket}{\rbrack\!\rbrack}
\providecommand{\euro}{{\text{€}}}
\providecommand{\no}{{\text{n°}}}
\providecommand{\interval}[2]{{[#1;#2]}}
\providecommand{\intervFO}[2]{{[#1;#2[}}
\providecommand{\dotfill}{\ldots}

\providecommand{\ds}{\displaystyle}
\providecommand{\textding}[1]{#1}
\providecommand{\Oij}{\(\mathrm{O};\vec{\imath},\vec{\jmath}\)}
\providecommand{\Oijk}{\(\mathrm{O};\vec{\imath},\vec{\jmath},\vec{k}\)}
\providecommand{\Ouv}{\(\mathrm{O};\vec{u},\vec{v}\)}
\providecommand{\intervFF}[2]{\ensuremath{[#1;#2]}}
\renewcommand{\euro}{\ifmmode\text{\texteuro}\else\texteuro\fi}
\DeclareUnicodeCharacter{00B0}{\ifmmode^{\circ}\else\textdegree\fi}
\DeclareUnicodeCharacter{20AC}{\ifmmode\text{\texteuro}\else\texteuro\fi}
% Énoncés : Éducation nationale. Transcription et corrigés : APMEP, Denis Vergès et contributeurs. Découpage et classement : CoopMaths, Sébastien Lozano, Éric Elter et contributeurs. Adaptation de la présentation : Maths Explorer.
% CC BY-SA : https://coopmaths.fr/www/about/
% https://coopmaths.fr/alea/static/eam/2026/tex/eam_specifique_2026_04_polynesie_1.tex
\begin{document}
\noindent\textit{APMEP / CoopMaths. CC BY-SA. Présentation adaptée par Maths Explorer.}\par\medskip
En 2025, une étude démographique indique qu'il y a eu \np{60000} naissances et \np{40000} décès dans une population donnée. Les spécialistes modélisent que le nombre de naissances diminuera de 2\,\% chaque année, tandis que le nombre de décès augmentera de 900 par an.

\medskip

Soit $n$ un entier naturel. On note :
\begin{itemize}[left=5mm]
	\item $u_{n}$ le nombre de naissances prévues l'année $2025+n$ ;
	\item $v_{n}$ le nombre de décès prévus l'année $2025+n$.
\end{itemize}

On a donc $u_{0}=\np{60000}$ et $v_{0}=\np{40000}$.

\begin{enumerate}
	\item
	\begin{enumerate}
		\item Déterminer le nombre de naissances prévues l'année 2026.
		\item Déterminer le nombre de décès prévus l'année 2026.
	\end{enumerate}

	\item
	\begin{enumerate}
		\item Justifier que pour tout entier naturel $n$, $u_{n+1}=0{,}98\,u_{n}$.
		\item La suite $(u_{n})$ est-elle arithmétique ou géométrique ? Préciser sa raison.
		\item Pour tout entier naturel $n$, exprimer $u_{n}$ en fonction de $n$.
	\end{enumerate}

	\item
	\begin{enumerate}
		\item Justifier que la suite $(v_{n})$ est arithmétique. Préciser sa raison.
		\item Pour tout entier naturel $n$, exprimer $v_{n}$ en fonction de $n$.
	\end{enumerate}

	\item La feuille de calcul ci-dessous donne les valeurs successives de $u_{n}$ et de $v_{n}$ (les valeurs sont évidemment arrondies à l'unité).

	\begin{center}
		\begin{tblr}[T]
			{hlines, vlines, columns={c}, column{1-2}={wd=1.5cm}, column{3-4}={wd=2.5cm},
				cells={font=\sffamily},
				row{1}={rowsep=0pt,ht=7mm, bg=gray9},row{2-Z}={m},
				column{1}={colsep=0pt,7mm,bg=gray9},
				cell{1}{1}={gray9,cmd={\tikz[baseline={(0,0.05)}]{\clip(0,0) rectangle (0.6,0.5) ;\fill[gray!50] (0,0)--(0.5,0.5)--(0.5,0);}}} }
			& A & B & C \\
			1  & n & $u_n$ & $v_n$ \\
			2  & 0 & \np{60000} & \np{40000} \\
			3  & 1 & \np{58800} & \np{40900} \\
			\vdots  & \vdots & \vdots & \vdots \\
			11 & 9 & \np{50025} & \np{48100} \\
			12 & 10 & \np{49024} & \np{49000} \\
			13 & 11 & \np{48044} & \np{49900} \\
			14 & 12 & \np{47083} & \np{50800} \\
			15 & 13 & \np{46141} & \np{51700} \\
		\end{tblr}
	\end{center}

	\begin{enumerate}
		\item Quelle formule doit-on saisir dans la cellule \textsf{C3} de sorte que, recopiée vers le bas, elle permette d'afficher les valeurs de la suite $(v_{n})$ ?
		\item Justifier, par un calcul, pourquoi \np{49000} est affiché dans la cellule \textsf{C12}.
	\end{enumerate}

	\item On appelle solde naturel annuel la différence entre le nombre de naissances et le nombre de décès de l'année. Ainsi le solde naturel en 2025 est \np{20000}.

	À l'aide de la feuille de calcul ci-dessus, déterminer à partir de quelle année le solde naturel de cette population devient négatif.
\end{enumerate}

\end{document}
